\documentclass[10pt,a4paper]{amsart}
\usepackage[margin=19mm]{geometry}
\usepackage{amsmath,amssymb,mathtools}
\usepackage[scr=rsfs,cal=euler]{mathalfa}
\usepackage{xfrac}
\usepackage[unicode,hidelinks,bookmarks=false]{hyperref}
\newcommand{\ie}{i.e.\ }
\newcommand{\cf}{cf.\ }
\newcommand{\resp}{respectively\ }
\newcommand{\Subsection}{\subsection}
\providecommand{\nobreakdash}{\nobreak\hbox{-}}
\setlength{\emergencystretch}{2em}
\multlinegap0pt
\usepackage{amssymb}
\usepackage{mathtools}
\usepackage{stackengine}
\usepackage{tikz}
\usepackage[normalem]{ulem}
\newtheorem{theorem}{Theorem}
\newtheorem{multicols}{Multicols}
\DeclareMathOperator{\Stab}{Stab}
\newcommand{\bT}{\mathsf{T}}
\DeclareMathOperator{\sgn}{sgn}
\title{Integrals of stable envelopes for cotangent bundles to Grassmannians}
\author{Matthew Crawford, Pavan Kartik, Reese Lance}
\date{Source: arXiv:2510.21573v2 (CC BY 4.0)}
\begin{document}
\maketitle

\noindent\emph{Source and licence.} Integrals of stable envelopes for cotangent bundles to Grassmannians. Version arXiv:2510.21573v2 (CC BY 4.0), distributed by arXiv under the Creative Commons Attribution 4.0 International licence, http://creativecommons.org/licenses/by/4.0/, as recorded in the arXiv licence metadata for this exact version and shown on the abstract page https://arxiv.org/abs/2510.21573v2. Full licence notice: \texttt{licenses/arxiv-2510.21573-CC-BY-4.0.txt}.\par\smallskip

\noindent\textit{Editorial scope.} Counted unit: the article's theorem lim-exists (source0.tex lines 1292--1294). The article's complete original proof of it is delivered with it (source0.tex lines 1295--1408, one \texttt{proof} environment of 114 lines). No mathematics is written, repaired or re-derived: the statement and the proof are the article's own lines, in the article's own notation, and no retained line is substituted or re-flowed.  No further source-local context is needed: every cross-reference of every retained line resolves inside the two delivered spans.  Nothing was omitted from any retained span, so the package carries no omission record.  No retained line contains a citation, so the wrapper carries no bibliography and no bibliography entry.  No retained line contains a figure environment, an image-inclusion command or drawing code, so no image is delivered and the package declares no asset dependency.  Editorial formatting only, in addition: an AMS \texttt{amsart} preamble that restates the article's own declarations and macros used by the retained lines, namely the environments \texttt{theorem}, \texttt{multicols} on separate counters, together with the standard packages those lines need.  The retained environments print in this document as Theorem 1, Multicols 1 in order of appearance, whereas the article prints the same items with the numbering of its own full document.  The complete native source of the article as distributed in the arXiv e-print for this exact version is bundled unchanged as \texttt{sources/187732/source0.tex}.
\begin{theorem}\label{lim-exists} The following limit exists for all $p_I\in X_{k,n}^\bT$:
\[\lim_{\mathbf{a}\to 0} \left( \sum_{p_J\in X_{k,n}^\bT}\frac{\Stab(p)}{e(T_{p_J}X_{k,n})}\right) \]
\end{theorem}

\begin{proof}~\\
Consider a fixed point $I=\{i_{1}<\dots<i_{k}\}$. The weight function is a rational function in $t_{1},\dots, t_{k}$ which is symmetric in $t_{1},\dots,t_{k}$ by definition, and the weight function of $p_{I}$ is given by
\[W(p_I)=\sum\limits_{\sigma \in S_k}\dfrac{\prod\limits_{r = 1}^k \prod\limits_{\alpha = 1}^{i_{r} - 1} (a_{\alpha} - t_{\sigma(r)}) \prod\limits_{\beta = i_r + 1}^n (t_{\sigma(r)} - a_{\beta} + \hbar) }{\prod\limits_{1 \leq \ell < m \leq k} (t_{\sigma(\ell)} - t_{\sigma(m)}) (t_{\sigma (\ell)} - t_{\sigma(m)} + \hbar)}\]
Note that 
\[
\sgn(\sigma)\cdot\prod\limits_{1\leq \ell<m\leq k}(t_{\ell}-t_{m})=\prod\limits_{1\leq \ell<m\leq k}(t_{\sigma(\ell)}-t_{\sigma(m)})\]
and
\[
\prod\limits_{1\leq l<m\leq k}(t_{\sigma(\ell)}-t_{\sigma(m)}+\hbar)(t_{\sigma(m)}-t_{\sigma(\ell)}+\hbar)=\prod\limits_{1\leq l<m\leq k}(t_{\ell}-t_{m}+\hbar)(t_{m}-t_{\ell}+\hbar)\] 
for any $\sigma\in S_{k}$.\\
Define \[
D:=\prod\limits_{1\leq\ell<m\leq k}(t_{\ell}-t_{m})(t_{\ell}-t_{m}+\hbar)(t_{m}-t_{\ell}+\hbar)
\]
The weight function associated to a fixed point may then be written as
\[
 W(p_{I})=\frac{N}{D}
\]
where $N,\ D$ are polynomials in the $a_{i}$, $t_{i}$ and $\hbar$. Since $W_{p_{I}}$ is a rational function in $t_{1},\dots, t_{k}$ and $D$ can also be thought of as a polynomial that is skew-symmetric in $t_{1},\dots,t_{k}$, we can conclude that numerator $N$ must also be skew symmetric in $t_{1},\dots,t_{k}$, which implies that the expression $V_{t}:=\prod\limits_{1\leq\ell<m\leq k}(t_{\ell}-t_{m})$ divides $N$. On specializing $t_{s}=a_{j_{s}}$ for any fixed point $J=\{j_{1}< \cdots < j_{k}\}$, we see that $N'_{J}:=\frac{N}{V_{t}}|_{t_{s}=a_{j_{s}}}$ and $D'_{J}:=\frac{D}{V_{t}}|_{t_{s}=a_{j_{s}}}$ are polynomials in the $a_{i}$ and $\hbar$. \\
For any ordered subset $X_{k,n}=\{x_{1},\dots,x_{p}\}$ define
\[V_{X_{k,n}}:=\prod\limits_{1\leq i<j\leq p}(x_{i}-x_{j})\]
For any fixed point $J=\{j_{1},\dots,j_{k}\}$, and permutation $\sigma \in S_k$, we define the following expressions:
\begin{multicols}{2}
\begin{itemize}
\item $V_{t}^{\hbar} := \prod\limits_{1 \leq \ell < m \leq k} (t_{\ell} - t_{m} + \hbar)$,
\item $\overline{V_{t}^{\hbar}} := \prod\limits_{1 \leq \ell < m \leq k} (t_{m} - t_{\ell} + \hbar)$,
\item $\overline{V_{J,\sigma}^{\hbar}} := \prod\limits_{1 \leq \ell < m \leq k} (a_{j_{\sigma(m)}} - a_{j_{\sigma(\ell)}} + \hbar)$,
\item $V_J := \prod\limits_{1 \leq \ell < m \leq k} (a_{j_\ell} - a_{j_m})$,
\item $V_J^{\hbar} :=  \prod\limits_{1 \leq \ell < m \leq k} (a_{j_\ell} - a_{j_m} + \hbar)$,
\item $\overline{V_J^{\hbar}} := \prod\limits_{1 \leq \ell < m \leq k} (a_{j_m} - a_{j_\ell} + \hbar)$,
\item $V_{J,J^c}^{\hbar} := \prod\limits_{v \not\in J} \prod\limits_{p = 1}^{k} (a_{j_p} - a_{v} + \hbar)$,
\item $D_t := V_t V_t^{\hbar} \overline{V_{t}^{\hbar}}$,
\item $D_J := V_J V_J^{\hbar} \overline{V_{J}^{\hbar}}$,
\item $D_J' := \frac{D_J}{V_J} = V_J^{\hbar}\overline{V_J^{\hbar}}$,
\end{itemize}
\end{multicols}
We also define \[N_{J,t} := \sum\limits_{\sigma \in S_k} \sgn(\sigma) \left(\prod\limits_{r = 1}^k \prod\limits_{\alpha = 1}^{i_{r} - 1} (a_{\alpha} - t_{\sigma(r)}) \prod\limits_{\beta = i_r + 1}^n (t_{\sigma(r)} - a_{\beta} + \hbar) \right) \overline{V_{J,\sigma}^{\hbar}}.\]
With these definitions, the restriction $W(p_I)|_{p_J}$ is
\[W(p_I)|_{p_J} = \left.\dfrac{N_{J,t}}{V_{t}V_{t}^{\hbar}\overline{V_{t}^{\hbar}}}\right|_{t_s = a_{j_s}}.\]
Note that the denominator before restriction is $D_t$.
Letting $N_{J}:=N_{J,t}|_{t_{s}=a_{j_{s}}}$, we can again write $W(p_I)|_{p_J}$ as \[W(p_I)|_{p_J} = \dfrac{N_J}{D_J}.\](Observe that $D_{t}$ is skew symmetric in the $t_{s}$'s (and specializes to $D_{J}$ when $t_{s}=a_{j_{s}}$). Since the weight function is symmetric, and the denominator is skew symmetric, the numerator $N_{J,t}$ must be skew symmetric as well, which implies that the Vandermonde determinant 
\[\prod\limits_{1 \leq \ell < m \leq k}(t_{\ell}-t_{m})\]
divides $N_{J,t}$, in particular the quotient is a polynomial in the $t_{s}'s$, with coefficients that are polynomials in terms of the $a_{s}'s$ and $\hbar$.\\\\
This implies that when we specialize to $t_{s}=a_{j_{s}}$, we obtain that
\[N'_{J}:=\frac{N_{J}}{V_{J}}\]
\[D'_{J} =\frac{D_{J}}{V_{J}}\]
are polynomials in the $a_{s}$'s and $\hbar$.\\\\


Hence we can write the integral as
\begin{align*}
\int_{X_{k,n}}\Stab(p_{I})&=\sum_{p_{J}\in X_{k,n}^{T}}\frac{\frac{N_{J}}{D_{J}}}{\prod\limits_{v\not \in J} \prod\limits_{p = 1}^{k}(a_v - a_{j_p}) (a_{j_p}-a_v+\hbar)}\\
&=\sum_{p_J\in X_{k,n}^T}\dfrac{\sum\limits_{\sigma \in S_k} \dfrac{\prod\limits_{r = 1}^k \prod\limits_{\alpha = 1}^{i_{r} - 1} \left(a_{\alpha} - a_{j_{\sigma(r)}}\right) \prod\limits_{\beta = i_r + 1}^n \left(a_{j_{\sigma(r)}} - a_{\beta} + \hbar\right) }{\prod\limits_{1 \leq \ell < m \leq k} \left(a_{j_{\sigma(\ell)}} - a_{j_{\sigma(m)}}\right) \left(a_{j_{\sigma(\ell)}} - a_{j_{\sigma(m)}} + \hbar\right)}}{\prod\limits_{v\not \in J} \prod\limits_{p = 1}^{k}\left(a_v - a_{j_p}\right) \left(a_{j_p}-a_v+\hbar\right)}
\end{align*}
\begin{align*}
&=\sum_{p_J\in X_{k,n}^T}\sum\limits_{\sigma \in S_k} \frac{\prod\limits_{r = 1}^k \prod\limits_{\alpha = 1}^{i_{r} - 1} \left(a_{\alpha} - a_{j_{\sigma(r)}}\right) \prod\limits_{\beta = i_r + 1}^n \left(a_{j_{\sigma(r)}} - a_{\beta} + \hbar\right) }{\prod\limits_{1 \leq \ell < m \leq k} \left(a_{j_{\sigma(\ell)}} - a_{j_{\sigma(m)}}\right) \left(a_{j_{\sigma(\ell)}} - a_{j_{\sigma(m)}} + \hbar\right)\prod\limits_{v\not \in J} \prod\limits_{p = 1}^{k}\left(a_v - a_{j_p}\right) \left(a_{j_p}-a_v+\hbar\right)}\\
&=\frac{1}{V}\sum_{p_{J}\in X_{k,n}^{T}}\frac{N_{J}V_{J^{c}}(-1)^{kn+\binom{k}{2}+|J|}}{D'_{J}V_{J,J^{c}}^{\hbar}}\\
&=\frac{1}{V}\sum_{p_{J}\in X_{k,n}^{T}}\frac{N'_{J}V_{J}V_{J^{c}}(-1)^{kn+\binom{k}{2}+|J|}}{D'_{J}V_{J,J^{c}}^{\hbar}},
\end{align*}
where $V_{J,J^c}^{\hbar}:= \prod_{v \not \in J}\prod_{p = 1}^{k} \left(a_{j_p} - a_v + \hbar\right)$. The factor of $(-1)^{kn+\binom{k}{2}+|J|}$ appears from rearranging the terms of $\prod\limits_{v\notin J}\prod\limits_{p=1}^{k}\left(a_{v}-a_{j_{p}}\right)$ so that $v<j_{p}$ for each term in the product.\newline
To show that the required limit exists, it is enough to show that the expression 
\begin{equation}\label{eq:thesum}\sum_{p_{J}\in X_{k,n}^{T}}\frac{N'_{J}V_{J}V_{J^{c}}(-1)^{k(n-k)+\binom{k}{2}+|J|}}{D'_{J}V_{J,J^{c}}^{\hbar}}\end{equation}
vanishes whenever we set $a_{x}=a_{y}=z$ for any pair $1\leq x<y\leq n$, since this would imply that $V$ divides this sum, which in turn implies that the limit exists. We consider the following cases, and fix $1\leq x<y\leq n$:\newline


\subsection{$J$ for which $x,\ y\in J$ or $x,y \not\in J$} 
Observe that the summands in \eqref{eq:thesum} corresponding to fixed points $p_J$ with both $x$ and $y$ in $J$ vanish, since $(a_x - a_y)$ will appear in $V_J$. Similarly, the summands in \eqref{eq:thesum} corresponding to fixed points $p_J$ for which neither $x$ nor $y$ are in $J$ vanish, since $(a_{x}-a_{y})$ will appear in $V_{J^{c}}$. \\

Note that we are now left with 2 cases: one for which $x\in J,\ y\notin J$, and one for which $x\notin J,\ y\in J$.\\
\subsection{$J$ contains exactly one of $x$ and $y$}
 Consider the summands corresponding to fixed points $p_{J}$ for which $y \in J$ and $x \notin J$. We can disregard such $p_{J}$ that are not in the attracting set of $p_{I}$, since the restriction $W(p_I)|_{p_J}$ is 0.\\


We will first show that if $p_{J_{y}^{x}}$ (where $J^{x}_{y}:=\left(J\backslash\{y\}\right)\cup\{x\}$) is not in the attracting set of $p_{I}$, then $W(p_{I})|_{p_{J}}=0$ 


\subsubsection{$(J\backslash\{y\})\cup\{x\}\notin Attr^{f}(I)$}
Since $y \in J$, there exists an index $1 \leq r \leq k$ such that $y = j_r$. Since $x \not\in J$, there exists an index $1 \leq p \leq r$ such that $j_{p-1} < x < j_p$ (the edge case of $p = 1$ is handled similarly). Then we have that:
\begin{align*}
J &= \{j_1 < \ldots < j_{p-1} < j_p < j_{p+1} < \ldots < j_{r-1} < \overbrace{j_r}^{y} < j_{r+1} < \ldots < j_k\}\\
J_y^x &= \{j_1 < \ldots < j_{p-1} < x< j_p < \ldots < j_{r-2} < j_{r-1} < j_{r+1} < \ldots < j_k\}\\
I &= \{i_1 < \ldots < i_{p-1} < i_p < i_{p+1} < \ldots < i_{r-1} < i_r < i_{r+1} < \ldots < i_k\}
\end{align*}
Since $J_y^x \not\in Att(I)$, some element of $J_{y}^{x}$ is less than the element of $I$ in the same position. Specifically, since the first $p-1$ elements of $J_{y}^{x}$ are the same as the first $p-1$ elements of $J$, and the last $k-r$ elements of $J_{y}^{x}$ are the same the last $k-r$ elements of $J$, the elements of $J_{y}^{x}$ that are less than the corresponding elements of $I$ must be in positions that are at least $p$ and  at most $r$. Therefore, either $x < i_p$ or $j_{b} < i_{b+1}$ for some $p \leq b \leq r-1$.\\ We note that for any permutation $\sigma \in S_k$, the summand in $N_J$ corresponding to $\sigma$ will be 0 if $x$ is less than the element of $I$ in the position corresponding to the position that $y$ was sent to under the action of $\sigma$, because there will be a factor of $a_x - a_y$ from the $\alpha$ sum which is 0 upon setting $a_x = a_y = z$.\\
Let $p \leq q \leq r$ be the largest index such that the element of $J_y^x$ in position $q$ is less than the element of $I$ in position $q$. Then we have:
\begin{align*}
J &= \{j_1 < \ldots < j_{p-1} < j_p < j_{p+1} < \ldots < j_{q-1} < j_q < j_{q+1} < \ldots < j_{r-1} < \overbrace{j_r}^{y} < j_{r+1} < \ldots < j_k\}\\
J_y^x &= \{j_1 < \ldots < j_{p-1} < x< j_p < \ldots < j_{q-2} < j_{q-1} < j_{q} < \ldots < j_{r-2} < j_{r-1} < j_{r+1} < \ldots < j_k\}\\
I &= \{i_1 < \ldots < i_{p-1} < i_p < i_{p+1} < \ldots < i_{q-1} < i_q < i_{q+1} < \ldots < i_{r-1} < i_r < i_{r+1} < \ldots < i_k\}
\end{align*}
In $J$, for the elements in position $b$ for $b \leq q-1$, we have that the greatest position that $j_b$ can be sent to under the action of a given permutation $\sigma \in S_k$ is $q-1$ since $j_b \leq j_{q-1} < i_q$ for all $b \leq q-1$. Therefore, if $\sigma \in S_k$ is a permutation that sends any of the first $q-1$ elements of $J$ to a position that is at least $q$, then, without loss of generality, the element in position $l$ of $J$ is sent to position $u$ in $\sigma(J)$ where $l < q \leq u 
\leq k$. But we know in the definition of $N_J$, in the $\alpha$ product, there will be a factor of 0 if there exists an element of $J$ that is permuted to a position where it is less than the corresponding element of $I$. We see that the $u^{\text{th}}$ element of $\sigma(J)$, which is $j_l$, which satisfies $j_l < j_{q-1} < i_q \leq i_u$, hence the summand corresponding to $\sigma$ is 0. Hence we see that it must be the case that $\sigma$ permutes the first $q-1$ elements of $J$ among themselves and $\sigma$ permutes the last $k-q+1$ elements of $J$ among themselves.\\
Since $y = j_r \geq j_q$, the position of $y$, under the action of $\sigma$, is again at least position $q$, that is $j_{\sigma^{-1}(r)} \geq j_q > i_q > x$. Also, $x < i_q < i_{\sigma^{-1}(r)}$ since $q < \sigma^{-1}(r)$.\\
Thus $y$, which is the element in position $\sigma^{-1}(r)$ of $\sigma(J)$ where $\sigma^{-1}(r) \geq q$, is at least $j_q$ since $r \geq q$, and by definition $j_q \geq i_q$, and $i_q > x$ by definition of $q$. We also note that since the elements of $I$ are increasing, $i_{\sigma^{-1}(r)} > i_q$ as $\sigma^{-1}(r) > q$, so $x<i_{\sigma^{-1}(r)}$, so the summand corresponding to $\sigma$ becomes 0 when setting $a_x = a_y = z$.\\\\
Note that now we are left with 2 types of summands: One type corresponds to fixed points $p_{J}$ in the attracting set of $p_{I}$ that contain $x$ and not $y$, and another type that corresponds to fixed points $p_{J}$ in the attracting set of $p_{I}$ that contain $y$ and not $x$. Our final step in this proof will be to show that the summands corresponding to $p_{J}$ and $p_{\hat{J}}$ add up to 0 when we set $a_{x}=a_{y}=z$, where $p_{J}$ is a fixed point in the attracting set of $p_{I}$ containing $y$ and not $x$, and $p_{\hat{J}}$ is a fixed point in the attracting set of $p_{I}$ obtained from $p_{J}$ by replacing $y$ with $x$. \\

\subsubsection{$\hat{J}:=(J\backslash\{y\})\cup\{x\}\in Attr^{f}(I)$}
Let $r$ denote the position of the smallest entry in $J$ so that $j_{r}>x$, and let $s$ denote the position of $y$ in $J$. \\
Consider the sum 
\begin{equation}\label{twosummands}
    \frac{N_{\hat{J}}V_{\hat{J}^{c}}(-1)^{k(n-k)+\binom{k}{2}+|\hat{J}|}}{D'_{\hat{J}}V_{\hat{J},\hat{J}^{c}}}+
    \frac{N_{J}V_{J^{c}}(-1)^{k(n-k)+\binom{k}{2}+|J|}}{D'_{J}V_{J,J^{c}}}
\end{equation}
Since $\hat{J}$ was obtained from $J$ by replacing $y$ with $x$, we have \[|\hat{J}|=|J|-y+x\]
The difference between $|J|$ and $|\hat{J}|$ contributes a sign of $(-1)^{y-x}$.\\
Observe that on setting $a_{x}=a_{y}=z$, the denominators become the same polynomial. The term $V_{\hat{J}^{c}}$ acquires a sign of $(-1)^{y-x-(r-s)}$ when setting $a_{x}=a_{y}=z$, i.e 
\[V_{\hat{J}^{c}}|_{a_{x}=a_{y}=z}=(-1)^{y-x-(r-s)}V_{J^{c}}|_{a_{x}=a_{y}=z}\]
Recall that $N_{J}$ and $N_{\hat{J}}$ are both defined as a sum over permutations $\sigma\in S_{k}$. Let $\tau\in S_{k}$ be the permutation so that $(\tau(\hat{J}))_{s}=x$, and $(\tau(\hat{J}))_{\ell}=j_{\ell}\ \forall \ell\neq s$. This permutation when applied to $\hat{J}$ shifts the elements to the ``correct position" as in $J$, which means that every element of $\hat{J}$ appears in the same place as in $J$, and $x$ appears in the same place as $y$. This permutation is a cycle with sign $(-1)^{s-r-1}$.\\\\
This means that we can now sum over permutations $\sigma\tau\in S_{k}$ (where $\sigma$ varies) for the $N_{\hat{J}}$ term in the summand that corresponds to $\hat{J}$ in \eqref{twosummands}. Doing this will allow us to factor out $\left(\frac{N_{J}V_{J^{c}}(-1)^{k(n-k)+\binom{k}{2}+|J|}}{D'_{J}V_{J,J^{c}}}\right)|_{a_{x}=a_{y}=z}$ from both summands in \eqref{twosummands}.\\  
From our previous observations we can see that the total sign that is obtained on the summand corresponding to $\hat{J}$ after setting $a_{x}=a_{y}=z$ in \eqref{twosummands} is \[(-1)^{-y+x+y-x-(r-s)+s-r-1}=-1\]
Which means that after setting $a_{x}=a_{y}=z$, \eqref{twosummands} can be written as
\[\left(\left(\frac{N_{J}V_{J^{c}}(-1)^{k(n-k)+\binom{k}{2}+|J|}}{D'_{J}V_{J,J^{c}}}\right)|_{a_{x}=a_{y}=z}\right)\times(-1+1)=0\]
Pairing up the other fixed points with this property (i.e $y\in J,\ x\notin J,\ (J\backslash\{y\})\cup\{x\}\in Attr^{f}(I)$) in a similar manner and following this procedure allows us to conclude that the sum is 0.\\\\
The results in the previous 3 subsections show that setting $a_{x}=a_{y}=z$ for arbitrary $1\leq x<y\leq n$ gives us 0 for $\sum_{p_{J}\in X_{k,n}^{T}}\frac{N_{J}V_{J^{c}}(-1)^{k(n-k)+\binom{k}{2}+|J|}}{D'_{J}V_{J,J^{c}}^{\hbar}}$, which in turn allows us to conclude that $V$ divides this sum. Hence the required limit exists.
\end{proof}

\end{document}
